\section{Noether normalization}
\label{section-Noether-normalization}

\noindent
In this section we prove variants of the Noether normalization lemma.
The key ingredient we will use is contained in the following two lemmas.

\begin{lemma}
\label{lemma-helper}
Let $n \in \mathbf{N}$.
Let $N$ be a finite nonempty
set of multi-indices $\nu = (\nu_1, \ldots, \nu_n)$.
Given $e = (e_1, \ldots, e_n)$ we set $e \cdot \nu = \sum e_i\nu_i$.
Then for $e_1 \gg e_2 \gg \ldots \gg e_{n-1} \gg e_n$ we have:
If $\nu, \nu' \in N$ then
$$
(e \cdot \nu = e \cdot \nu')
\Leftrightarrow
(\nu = \nu')
$$
\end{lemma}

\begin{proof}
Say $N = \{\nu_j\}$ with $\nu_j = (\nu_{j1}, \ldots, \nu_{jn})$.
Let $A_i = \max_j \nu_{ji} - \min_j \nu_{ji}$. If for each $i$ we have
$e_{i - 1} > A_ie_i + A_{i + 1}e_{i + 1} + \ldots + A_ne_n$ then
the lemma holds. For suppose that $e \cdot (\nu - \nu') = 0$. Then for
$n \ge 2$,
$$
e_1(\nu_1 - \nu'_1) = \sum\nolimits_{i = 2}^n e_i(\nu'_i - \nu_i).
$$
We may assume that $(\nu_1 - \nu'_1) \ge 0$. If $(\nu_1 - \nu'_1) > 0$, then
$$
e_1(\nu_1 - \nu'_1) \ge e_1 >
A_2e_2 + \ldots + A_ne_n \ge
\sum\nolimits_{i = 2}^n e_i|\nu'_i - \nu_i| \ge
\sum\nolimits_{i = 2}^n e_i(\nu'_i - \nu_i).
$$
This contradiction implies that $\nu'_1 = \nu_1$. By
induction, $\nu'_i = \nu_i$ for $2 \le i \le n$.
\end{proof}

\begin{lemma}
\label{lemma-helper-polynomial}
Let $R$ be a ring. Let $g \in R[x_1, \ldots, x_n]$ be an element
which is nonconstant, i.e., $g \not \in R$.
For $e_1 \gg e_2 \gg \ldots \gg e_{n-1} \gg e_n = 1$ the polynomial
$$
g(x_1 + x_n^{e_1}, x_2 + x_n^{e_2}, \ldots, x_{n - 1} + x_n^{e_{n - 1}}, x_n)
=
ax_n^d + \text{lower order terms in }x_n
$$
where $d > 0$ and $a \in R$ is one of the nonzero coefficients of $g$.
\end{lemma}

\begin{proof}
Write $g = \sum_{\nu \in N} a_\nu x^\nu$ with $a_\nu \in R$ not zero.
Here $N$ is a finite set of multi-indices as in
Lemma \ref{lemma-helper}
and $x^\nu = x_1^{\nu_1} \ldots x_n^{\nu_n}$.
Note that the leading term in
$$
(x_1 + x_n^{e_1})^{\nu_1} \ldots (x_{n-1} + x_n^{e_{n-1}})^{\nu_{n-1}}
x_n^{\nu_n}
\quad\text{is}\quad
x_n^{e_1\nu_1 + \ldots + e_{n-1}\nu_{n-1} + \nu_n}.
$$
Hence the lemma follows from
Lemma \ref{lemma-helper}
which guarantees that there is exactly one nonzero term $a_\nu x^\nu$ of $g$
which gives rise to the leading term of
$g(x_1 + x_n^{e_1}, x_2 + x_n^{e_2}, \ldots, x_{n - 1} + x_n^{e_{n - 1}},
x_n)$, i.e., $a = a_\nu$ for the unique $\nu \in N$ such that
$e \cdot \nu$ is maximal.
\end{proof}

\begin{lemma}
\label{lemma-one-relation}
Let $k$ be a field.
Let $S = k[x_1, \ldots, x_n]/I$ for some proper ideal $I$.
If $I \not = 0$, then there exist $y_1, \ldots, y_{n-1} \in k[x_1, \ldots, x_n]$
such that $S$ is finite over $k[y_1, \ldots, y_{n-1}]$. Moreover we may
choose $y_i$ to be in the $\mathbf{Z}$-subalgebra of $k[x_1, \ldots, x_n]$
generated by $x_1, \ldots, x_n$.
\end{lemma}

\begin{proof}
Pick $f \in I$, $f\not = 0$. It suffices to show the lemma
for $k[x_1, \ldots, x_n]/(f)$ since $S$ is a quotient of that ring.
We will take $y_i = x_i - x_n^{e_i}$, $i = 1, \ldots, n-1$
for suitable integers $e_i$. When does this work? It suffices
to show that $\overline{x_n} \in k[x_1, \ldots, x_n]/(f)$
is integral over the ring $k[y_1, \ldots, y_{n-1}]$. The
equation for $\overline{x_n}$ over this ring is
$$
f(y_1 + x_n^{e_1}, \ldots, y_{n-1} + x_n^{e_{n-1}}, x_n) = 0.
$$
Hence we are done if we can show there exists integers $e_i$ such
that the leading coefficient with respect to $x_n$ of the equation
above is a nonzero element of $k$. This can be achieved for example
by choosing $e_1 \gg e_2 \gg \ldots \gg e_{n-1}$, see
Lemma \ref{lemma-helper-polynomial}.
\end{proof}

\begin{lemma}
\label{lemma-Noether-normalization}
\begin{slogan}
Noether normalization
\end{slogan}
Let $k$ be a field. Let $S = k[x_1, \ldots, x_n]/I$ for some ideal $I$.
If $I \neq (1)$, there exist $r\geq 0$, and
$y_1, \ldots, y_r \in k[x_1, \ldots, x_n]$
such that (a) the map $k[y_1, \ldots, y_r] \to S$ is injective
where the source is the polynomial ring on $y_1, \ldots, y_r$,
and (b) the map $k[y_1, \ldots, y_r] \to S$ is finite.
In this case the integer $r$ is the dimension of $S$.
Moreover we may choose $y_i$ to be in the
$\mathbf{Z}$-subalgebra of $k[x_1, \ldots, x_n]$
generated by $x_1, \ldots, x_n$.
\end{lemma}

\begin{proof}
By induction on $n$, with $n = 0$ being trivial.
If $I = 0$, then take $r = n$ and $y_i = x_i$.
If $I \not = 0$, then choose $y_1, \ldots, y_{n-1}$
as in Lemma \ref{lemma-one-relation}. Let
$S' \subset S$ be the subring generated by
the images of the $y_i$. By induction we can
choose $r$ and $z_1, \ldots, z_r \in k[y_1, \ldots, y_{n-1}]$
such that (a), (b) hold for $k[z_1, \ldots, z_r]
\to S'$. Since $S' \to S$ is injective and finite
we see (a), (b) hold for $k[z_1, \ldots, z_r]
\to S$. The last assertion follows from Lemma
\ref{lemma-integral-sub-dim-equal}.
\end{proof}

\begin{lemma}
\label{lemma-Noether-normalization-at-point}
Let $k$ be a field.
Let $S$ be a finite type $k$ algebra and denote $X = \Spec(S)$.
Let $\mathfrak q$ be a prime of $S$, and let $x \in X$ be the
corresponding point. There exists a $g \in S$, $g \not \in \mathfrak q$
such that $\dim(S_g) = \dim_x(X) =: d$ and such that
there exists a finite injective map $k[y_1, \ldots, y_d] \to S_g$.
\end{lemma}

\begin{proof}
Note that by definition $\dim_x(X)$ is the minimum
of the dimensions of $S_g$ for $g \in S$, $g \not \in \mathfrak q$, i.e.,
the minimum is attained. Thus the lemma follows from
Lemma \ref{lemma-Noether-normalization}.
\end{proof}

\begin{lemma}
\label{lemma-refined-Noether-normalization}
Let $k$ be a field. Let $\mathfrak q \subset k[x_1, \ldots, x_n]$
be a prime ideal. Set $r = \text{trdeg}_k\ \kappa(\mathfrak q)$.
Then there exists a finite ring map
$\varphi : k[y_1, \ldots, y_n] \to k[x_1, \ldots, x_n]$ such
that $\varphi^{-1}(\mathfrak q) = (y_{r + 1}, \ldots, y_n)$.
\end{lemma}

\begin{proof}
By induction on $n$. The case $n = 0$ is clear. Assume $n > 0$.
If $r = n$, then $\mathfrak q = (0)$ and the result is clear.
Choose a nonzero $f \in \mathfrak q$. Of course $f$ is nonconstant.
After applying an automorphism of the form
$$
k[x_1, \ldots, x_n] \longrightarrow k[x_1, \ldots, x_n],
\quad
x_n \mapsto x_n,
\quad
x_i \mapsto x_i + x_n^{e_i}\ (i < n)
$$
we may assume that $f$ is monic in $x_n$ over $k[x_1, \ldots, x_n]$, see
Lemma \ref{lemma-helper-polynomial}. Hence the ring map
$$
k[y_1, \ldots, y_n] \longrightarrow k[x_1, \ldots, x_n],
\quad
y_n \mapsto f,
\quad
y_i \mapsto x_i\ (i < n)
$$
is finite. Moreover $y_n \in \mathfrak q \cap k[y_1, \ldots, y_n]$ by
construction. Thus
$\mathfrak q \cap k[y_1, \ldots, y_n] = \mathfrak pk[y_1, \ldots, y_n] + (y_n)$
where $\mathfrak p \subset k[y_1, \ldots, y_{n - 1}]$ is a prime ideal.
Note that $\kappa(\mathfrak p) \subset \kappa(\mathfrak q)$ is finite, and
hence $r = \text{trdeg}_k\ \kappa(\mathfrak p)$.
Apply the induction hypothesis to the pair
$(k[y_1, \ldots, y_{n - 1}], \mathfrak p)$ and we obtain a finite ring map
$k[z_1, \ldots, z_{n - 1}] \to k[y_1, \ldots, y_{n - 1}]$ such that
$\mathfrak p \cap k[z_1, \ldots, z_{n - 1}] = (z_{r + 1}, \ldots, z_{n - 1})$.
We extend the ring map
$k[z_1, \ldots, z_{n - 1}] \to k[y_1, \ldots, y_{n - 1}]$
to a ring map
$k[z_1, \ldots, z_n] \to k[y_1, \ldots, y_n]$
by mapping $z_n$ to $y_n$.
The composition of the ring maps
$$
k[z_1, \ldots, z_n] \to k[y_1, \ldots, y_n] \to k[x_1, \ldots, x_n]
$$
solves the problem.
\end{proof}

\begin{lemma}
\label{lemma-Noether-normalization-over-a-domain}
Let $R \to S$ be an injective finite type ring map. Assume $R$ is a domain.
Then there exists an integer $d$ and a factorization
$$
R \to R[y_1, \ldots, y_d] \to S' \to S
$$
by injective maps such that $S'$ is finite over $R[y_1, \ldots, y_d]$
and such that $S'_f \cong S_f$ for some nonzero $f \in R$.
\end{lemma}

\begin{proof}
Pick $x_1, \ldots, x_n \in S$ which generate $S$ over $R$.
Let $K$ be the fraction field of $R$ and $S_K = S \otimes_R K$.
By Lemma \ref{lemma-Noether-normalization}
we can find $y_1, \ldots, y_d \in S$ such that $K[y_1, \ldots, y_d] \to S_K$
is a finite injective map. Note that $y_i \in S$ because we may pick the
$y_j$ in the $\mathbf{Z}$-algebra generated by $x_1, \ldots, x_n$.
As a finite ring map is integral (see
Lemma \ref{lemma-finite-is-integral})
we can find monic $P_i \in K[y_1, \ldots, y_d][T]$ such that
$P_i(x_i) = 0$ in $S_K$. Let $f \in R$ be a nonzero element such that
$fP_i \in R[y_1, \ldots, y_d][T]$ for all $i$. Then $fP_i(x_i)$
maps to zero in $S_K$. Hence after replacing $f$ by another
nonzero element of $R$ we may also assume $fP_i(x_i)$ is zero in $S$.
Set $x_i' = fx_i$ and let $S' \subset S$ be the $R$-subalgebra generated by
$y_1, \ldots, y_d$ and $x'_1, \ldots, x'_n$. Note that $x'_i$ is integral over
$R[y_1, \ldots, y_d]$ as we have $Q_i(x_i') = 0$ where
$Q_i = f^{\deg_T(P_i)}P_i(T/f)$ which is a monic polynomial in
$T$ with coefficients in $R[y_1, \ldots, y_d]$ by our choice of $f$.
Hence $R[y_1, \ldots, y_d] \subset S'$ is finite by
Lemma \ref{lemma-characterize-finite-in-terms-of-integral}.
Since $S' \subset S$ we have $S'_f \subset S_f$ (localization is exact).
On the other hand, the elements $x_i = x'_i/f$ in $S'_f$ generate $S_f$
over $R_f$ and hence $S'_f \to S_f$ is surjective. Whence
$S'_f \cong S_f$ and we win.
\end{proof}





